\documentclass[manuscript,11pt,screen,nonacm]{acmart}
\geometry{letterpaper,margin=1in,heightrounded=false,headheight=14pt,footskip=24pt}
\fancypagestyle{submission}{\fancyhf{}\fancyfoot[C]{\normalfont\normalsize\thepage}
  \renewcommand{\headrulewidth}{0pt}\renewcommand{\footrulewidth}{0pt}}
\setcopyright{none}
\settopmatter{printacmref=false,printccs=false,printfolios=true}
\acmDOI{}\acmISBN{}\acmVolume{}\acmNumber{}
% This is an unsubmitted manuscript. Keep the official class unchanged,
% and disable hyperxmp's automatic ACM publisher/journal inference.
\makeatletter\let\hyxmp@parse@acmart\relax\makeatother
\hypersetup{keeppdfinfo=true,colorlinks=true,linkcolor=blue,citecolor=blue,urlcolor=blue,
 pdftitle={Certified Expected-Cost Optimization for Robust Ski Rental},
 pdfauthor={Zijian Zeng, Jiayin Cao, Weinan Dai, Rui Wang, Fei Ding},
 pdfpublisher={},pdfpubtype={},pdfmetalang={en},
 pdfcontactemail={1002266693@ucsiuniversity.edu.my, 2795820153@qq.com, 21725058@zju.edu.cn, 2002007@nut.edu.cn, dingfei@email.ncu.edu.cn},
 pdfcontactaddress={UCSI University; Ningbo University; Zhejiang University; Nanchang Institute of Technology; Alibaba Group},
 pdfsubject={Robust randomized ski rental with distributional advice},
 pdfkeywords={ski rental, distributional advice, robustness, exact optimization}}
\renewcommand{\bibliofont}{\normalfont\normalsize}
\captionsetup{font=normalsize,labelfont={bf,normalsize},textfont=normalsize}
\newtheorem{theorem}{Theorem}
\newtheorem{lemma}[theorem]{Lemma}
\newtheorem{proposition}[theorem]{Proposition}
\DeclareMathOperator{\supp}{supp}
\newcommand{\E}{\mathbb E}
\newcommand{\Q}{\mathcal F}
\newcommand{\Pp}{\mathcal P}
\usepackage{tikz}
\usetikzlibrary{arrows.meta}
\usepackage{algorithm}
\usepackage{needspace}
\title{Certified Expected-Cost Optimization for Robust Ski Rental}
\author{Zijian Zeng}
\affiliation{\institution{UCSI University}\city{}\country{Malaysia}}
\email{1002266693@ucsiuniversity.edu.my}
\author{Jiayin Cao}
\affiliation{\institution{Ningbo University}\city{}\country{China}}
\email{2795820153@qq.com}
\author{Weinan Dai}
\affiliation{\institution{Zhejiang University}\city{}\country{China}}
\email{21725058@zju.edu.cn}
\author{Rui Wang}
\affiliation{\institution{Nanchang Institute of Technology}\city{}\country{China}}
\email{2002007@nut.edu.cn}
\author{Fei Ding}
\affiliation{\institution{Alibaba Group}\city{}\country{China}}
\email{dingfei@email.ncu.edu.cn}
\authorsaddresses{\normalsize\begin{tabular}{@{}ll@{}}
Zijian Zeng & \href{mailto:1002266693@ucsiuniversity.edu.my}{1002266693@ucsiuniversity.edu.my}\\
Jiayin Cao & \href{mailto:2795820153@qq.com}{2795820153@qq.com}\\
Weinan Dai & \href{mailto:21725058@zju.edu.cn}{21725058@zju.edu.cn}\\
Rui Wang & \href{mailto:2002007@nut.edu.cn}{2002007@nut.edu.cn}\\
Fei Ding & \href{mailto:dingfei@email.ncu.edu.cn}{dingfei@email.ncu.edu.cn}
\end{tabular}\spacefactor=3000}
\date{}
\begin{document}
\singlespacing
\begin{abstract}
\normalsize
We study discrete ski rental with integer purchase cost $b\ge2$, a finite
rational forecast of season length, and a prescribed expected competitive
ratio at every actual integer horizon. We give an exact algorithm that
minimizes forecast expected cost and returns matching primal and dual
certificates on the unbounded horizon and buy-day domains. It uses a
triangular prefix optimization, normalizes the policy by increasing a dual
price, and recovers the robustness constraint through one moment price.
An optimum needs at most $b+1$ buy-days, including at most two days at or
after $b$. For the uniform forecast on $\{1,\ldots,b\}$, we derive all $b$
affine decreasing pieces of the value frontier and its constant continuation,
prove uniqueness up to robustness two, and show that both support bounds
are sharp. The underlying capped-geometric policy is credited to prior
one-hot-forecast work. An exact example separates expected-cost optimization
from a published minimum-support-level inference; a pinned software
reproduction gives a third, distinct value. The algorithm has polynomial
bit time in numerical $b$ and the binary input length.
\end{abstract}
\hypersetup{pageanchor=false}
\maketitle
\hypersetup{pdftitle={Certified Expected-Cost Optimization for Robust Ski Rental},
 pdfauthor={Zijian Zeng, Jiayin Cao, Weinan Dai, Rui Wang, Fei Ding},
 pdfpublisher={},pdfpubtype={},pdfsubject={Robust randomized ski rental with distributional advice},
 pdfkeywords={ski rental, distributional advice, robustness, exact optimization}}
\enlargethispage{-12pt}\thispagestyle{empty}\clearpage
\setcounter{page}{1}\hypersetup{pageanchor=true}\pagestyle{submission}

\section{Introduction}\label{sec:intro}
A forecast can identify the season lengths that matter most, while a
competitive guarantee protects against seasons the forecast misses. The
resulting design question is concrete: among all policies meeting a chosen
worst-case guarantee, which one has the smallest forecast expected cost?
For randomized ski rental, balancing these requirements involves the whole
buying-day distribution, rather than a single threshold.

Renting costs one per day and buying costs an integer $b\ge2$.
A policy draws a positive integer buy-day $T$ independently of the actual
season length $x$ and pays
\[
 C_t(x)=\begin{cases}x,&t>x,\\ b+t-1,&t\le x.\end{cases}
\]
The input $p$ is a finite rational probability law on positive integer
horizons. Write $g(t)=\sum_xp_x C_t(x)$ and $J_p(f)=\sum_tg(t)f_t$.
For rational $R\ge\rho_b:=b^b/(b^b-(b-1)^b)$, we minimize $J_p(f)$ subject to
\begin{equation}\label{eq:robust}
 \E_f C_T(x)\le R\min(x,b)\qquad\text{for every integer }x\ge1.
\end{equation}
Thus the forecast determines the objective, while actual horizons remain
unbounded.

A small example explains why the choice of objective matters. At $b=3$,
$R=3/2$, and $p=(2/5,3/5)$ on horizons one and two, the global optimum is
$39/20$. Restricting the policy to the smallest feasible level of its raw
support costs instead gives $11/5$. The publisher-linked software, in six
specified calls at one pinned version, returns a law costing $12/5$.
Section~\ref{sec:separation} proves the first two statements and records the
last as a finite software observation. This separates the optimality
inference in Kim and Fan's published Appendix D.2~\cite{pmlr-v306-kim26l}
from both ideal support feasibility and literal implementation behavior.

\paragraph{Contributions and the LP comparison.}
We give an explicit certifying replacement for that inference. After a
finite support reduction, a backward elimination solves the prefix problem;
a normalization price activates constraints in a nested order; and a moment
price recovers the prescribed robustness. The algorithm returns a global
lower-bound witness as well as an optimal policy. For a uniform forecast on
$\{1,\ldots,b\}$, we further obtain the complete expected-cost frontier,
strict-dual uniqueness, and simultaneous necessity of $b+1$ support days and
two days at or after $b$.

A generic rational LP already gives polynomial bit time after the finite
reduction. Our contribution is the explicit construction and its structural
consequences, with the same numerical-$b$ runtime model. In particular, the
moment row contains $t-1$, which may be large when a forecast horizon is
binary encoded; only the prefix block has coefficients bounded by $b$.
Tardos's matrix-dependent arithmetic bound~\cite{Tardos_1986} therefore
does not make this whole matrix small. Arithmetic-step bounds and bit time
are distinct, and we claim neither a strongly polynomial algorithm nor an
asymptotic improvement over general LP.

\paragraph{Context and closest results.}
Classical rent-or-buy analysis grew from snoopy caching~\cite{Karlin1988};
randomization gives the limiting ratio $e/(e-1)$~\cite{Karlin1994}.
Point-prediction robustness--consistency tradeoffs and LP methods appear in
Purohit, Svitkina and Kumar~\cite{NEURIPS2018_73a427ba} and Wei and
Zhang~\cite{NEURIPS2020_5bd844f1}. Li, Christianson and Li~\cite{Li_2026}
develop prediction-specific Pareto optimization, equalizing distributions
and primal--dual constructions. Their inspected preprint's randomized
strong-optimality statements are asymptotic in $b$.

The nearest model is Kim and Fan~\cite{pmlr-v306-kim26l}. We use their
Lemma 5.1 robustness reduction and credit their geometric and gap
recurrences. The uniform policy below coincides with their Theorem C.1(ii) in Appendix C.4
one-hot-$b$ optimizer when $R-1\le1-1/b$, and is another one-hot-$b$ optimum
above that cut. The uniform contribution is its expected-objective frontier
and uniqueness analysis, with this underlying policy structure retained.

Distributional learning~\cite{pmlr-v139-diakonikolas21a} combines a learned
policy with a robust baseline and does provide per-instance guarantees;
we optimize exactly at each supplied $p,R$. Besbes, Ma and
Mouchtaki~\cite{Besbes_2025} use heterogeneous data and known-radius
shift to study additive regret. Cui and Dinitz~\cite{pmlr-v306-cui26a}
consider unknown Wasserstein prediction error and additionally interpolate
with robustness. Their benchmark is a policy knowing the true distribution.
Khanafer, Kodialam and Puttaswamy~\cite{Khanafer_2013} use horizon-moment
information; our moment is instead a property of the buying-day law.
Tail-risk constraints~\cite{Dinitz_2024,10.1007/978-3-032-06706-7_6}
bound the probability of a large realized ratio, while here an expected-ratio
cap is fixed at each horizon. Hartline, Johnsen and
Shah~\cite{hartline_et_al:LIPIcs.ITCS.2026.75} study finite binary-weather
sequences, with a distribution-knowing benchmark in their prior-independent
analysis and a separate subgame-optimality criterion.

Set $a=b-1$, $r=R-1$, $K=rb$, and
\begin{equation}\label{eq:S}
 S=\{1,\ldots,b\}\cup\{d+1:p_d>0,\ d\ge b\}.
\end{equation}
\begin{theorem}\label{thm:main}
For every input above, the explicit algorithm in Sections~\ref{sec:oracle}
and~\ref{sec:price} returns an expected-cost minimizing robust law $f$,
supported on $S$, with at most two support days $t\ge b$ and at most $b+1$
support days in all. It returns rational $\lambda\ge0$, $y_x\ge0$ for
$1\le x<b$, and $\alpha\in\mathbb Q$ such that
\begin{align}
 g(t)+\lambda(t-1)+\sum_{x=t}^{b-1}(b+t-1-x)y_x&\ge\alpha
           &&(t\ge1),\label{eq:certificate}\\
 J_p(f)&=\alpha-r\sum_{x=1}^{b-1}xy_x-\lambda K.\label{eq:equality}
\end{align}
Empty sums are zero. Let $q_R$ be the denominator of $R$, $q_p$ the least
common multiple of the probability denominators, and
\begin{equation}\label{eq:bound}
 D=b!b^b,\qquad H=q_pq_RD,\qquad L=b+\max S,\qquad Z=2LH^4.
\end{equation}
After computing $g$ on $S$, the algorithm uses $O(\log Z)$ calls to an
$O(b^2+|S|)$ rational-operation oracle. All intermediate bit lengths and
the total bit time are polynomial in numerical $b$ and binary input length.
\end{theorem}

Figure~\ref{fig:pipeline} shows how the two prices separate normalization
from robustness. The resulting certificates can be checked using the
prefix constraints, one moment, and finitely many inequalities on $S$.
\begin{figure}[ht]
\centering
\begin{tikzpicture}[>=Stealth,every node/.style={font=\normalsize},
 box/.style={draw,rounded corners,align=center,inner sep=5pt}]
\node[box] (g) at (0,0) {$g$ on $S$};
\node[box] (pre) at (2.7,0) {Prefix elimination\\normalization price $\alpha$};
\node[box] (price) at (6.3,0) {Moment price $\lambda$\\exact bisection};
\node[box] (mix) at (10.2,0) {Policy mixture\\primal/dual certificate};
\draw[->] (g)--(pre);\draw[->] (pre)--(price);\draw[->] (price)--(mix);
\end{tikzpicture}
\caption{The normalized prefix oracle is called at adjusted costs
$g(t)+\lambda(t-1)$; the final mixture enforces the moment cap.}
\Description{Forecast costs enter the normalized prefix oracle, then moment-price bisection, then an optimal mixture and its dual certificate.}
\label{fig:pipeline}
\end{figure}

\section{Expected cost and raw support levels}\label{sec:separation}
The raw support objective is $H(f)=\sup_{t\in\supp f}g(t)$; on the finite
set $S$ this is a maximum. Let $h_*:=\inf_{f\text{ robust}}H(f)$.
Minimizing $H$ controls every supported buy-day, whereas minimizing $J_p$
allows a costly day if its probability is small. These objectives can
therefore choose different supports.

The published Appendix D.2, printed page 30, states:
\begin{quote}
``Therefore the optimal objective value equals the minimum feasible h,
which can be found by binary search.''
\end{quote}
Here $h$ is the raw support level, not a dual-adjusted expected objective.
The following example refutes that inference in the published model
\cite{pmlr-v306-kim26l}; it also applies to the corresponding arXiv v1 passage.

\Needspace{12\baselineskip}
\begin{proposition}\label{prop:separation}
For $b=3$, $R=3/2$, $p_1=2/5$, $p_2=3/5$, the global expected-cost
minimum is $39/20$, whereas $h_*=14/5$ and the unique feasible law at this
minimum support level has expected cost $11/5$.
\end{proposition}
\begin{proof}
Direct summation gives $g=(3,14/5,8/5,8/5,\ldots)$.
Robustness implies $\mu(f):=\E(T-1)\le3/2$, as proved in Section~\ref{sec:reduction}.
For all $t\ge1$, $g(t)+(7/10)(t-1)\ge3$, with equality only at 1 and 3.
Hence $J_p(f)\ge39/20$. The law $(f_1,f_3)=(1/4,3/4)$ attains this bound.
Its costs at actual horizons 1, 2 and $x\ge3$ are $3/2,9/4,9/2$, so it is robust.
Below level $14/5$, all allowed days are at least 3 and have moment at least 2.
At level $14/5$, $f_1=0$. The horizon-two constraint gives $2f_2\le1$,
while $\mu\ge2-f_2$. The cap forces $f_2=1/2$ and all remaining mass at
3. This robust law has $J_p=11/5$, proving the assertions.
\end{proof}

\begin{table}[ht]
\centering\normalsize
\caption{Three distinct quantities for Proposition~\ref{prop:separation}.
The last column is the raw support objective $H(f)$ whose minimum is $h_*$.
The software row is a finite observation; the other two rows are exact optima.}
\label{tab:separation}
\begin{tabular}{lcccc}
\toprule
Meaning & $(f_1,f_2,f_3)$ & $J_p$ & $\mu$ & $H(f)$\\
\midrule
Global expected-cost optimum & $(1/4,0,3/4)$ & $39/20$ & $3/2$ & $3$\\
Minimum support-level law & $(0,1/2,1/2)$ & $11/5$ & $3/2$ & $14/5$\\
Pinned software output & $(1/4,3/8,3/8)$ & $12/5$ & $9/8$ & $3$\\
\bottomrule
\end{tabular}
\end{table}

\paragraph{Published software observation.}
We executed the unchanged \texttt{SkiRentalOptimizer} class in both
optimizer files supporting Figure 1 and Table 2 in the publisher-linked
repository. The exact pinned version, file hashes and calls are supplied
in the reproduction supplement.
For constructor arguments $(3,1.5,[1,2],[.4,.6])$ and tolerances
$10^{-5},10^{-9},10^{-12}$, all six calls returned the last law in
Table~\ref{tab:separation}, at levels approximately 3. Both classes rejected
$14/5$ and accepted 3 in the documented feasibility calls. At level 3 the
construction temporarily assigns day-three mass $9/16$, reaching internal
mass $19/16$, before its existing correction leaves $3/8$. Evaluating the
returned law in the original model gives moment $9/8$ and horizon costs
$3/2,3,33/8$ at $1,2,x\ge3$, respectively. Thus the policy is robust and
costs $12/5$. These calls concern the pinned files and tolerances; the
support-level counterexample holds independently of their numerical behavior.

\paragraph{A different normalization changes the optimum.}
Dividing $J_p$ by the fixed $\E_p\min(X,b)$ preserves all optimizers.
In contrast, for $b=3,R=3/2,p_1=p_2=1/2$, define
$k(t)=\E_p[C_t(X)/\min(X,b)]$. Here
$g=(3,5/2,3/2,\ldots)$ and $k=(9/4,3/2,1,\ldots)$.
The global inequalities
\[
 g(t)+\tfrac34(t-1)\ge3,\qquad
 k(t)+\tfrac12(t-1)\ge2
\]
are equalities only on $\{1,3\}$ and $\{2,3\}$, respectively.
Together with $\mu\le3/2$, they prove unique optimal laws
$(f_1,f_3)=(1/4,3/4)$ for $J_p$, value $15/8$, and
$(f_2,f_3)=(1/2,1/2)$ for $\E_f k(T)$, value $5/4$.
Both laws satisfy the prefix rows. Their cross-objective values are
$21/16$ and $2$, respectively.

\section{Robustness and finite support}\label{sec:reduction}
Expanding the expected cost gives
\[
 \E C_T(x)=x+\sum_{t\le x}(t-1)f_t+(b-x)\sum_{t\le x}f_t.
\]
For each fixed $t$, the cost is nondecreasing in $x$ and converges to
$b+t-1$. Monotone convergence shows that~\eqref{eq:robust} implies
$\mu(f):=\E(T-1)\le K$. Conversely, this moment bound and the $a$ prefix
inequalities
\begin{equation}\label{eq:prefix}
 \sum_{t\le x}(a+t-x)f_t\le rx,\qquad 1\le x\le a,
\end{equation}
imply~\eqref{eq:robust}: the displayed inequalities cover $x<b$, and
$\E C_T(x)\le b+\mu(f)\le Rb$ covers $x\ge b$.
This is the robustness reduction of~\cite[Lemma 5.1]{pmlr-v306-kim26l}.

For $t\ge b$, let $s$ be the largest member of $S\cap[b,\infty)$ not
exceeding $t$. No positive forecast atom lies in $[s,t-1]$, so
\begin{equation}\label{eq:dominance}
 g(t)-g(s)=(t-s)\Pr_{X\sim p}[X\ge t]\ge0.
\end{equation}
Moving mass from $t$ to $s$ decreases its moment, preserves every prefix
constraint, and does not increase its objective. The same dominance holds
for $g(t)+\lambda(t-1)$ when $\lambda\ge0$.
This elementary reduction extends the point-prediction support principle
of~\cite[Appendix E.1, arXiv v1]{li2025predictionspecificdesignlearningaugmentedalgorithms}.

The finite robust polytope on $S$ is compact and nonempty. Indeed, the law
\[
 f_t=\frac{b^{t-1}(b-1)^{b-t}}{b^b-(b-1)^b},\quad 1\le t\le b,
\]
sums to one, satisfies~\eqref{eq:prefix} with equality at $R=\rho_b$,
and has moment $b(\rho_b-1)$. These identities follow by geometric summation,
or by the tight-constraint recurrence below. It remains feasible for larger
$R$. Hence a finite-support optimum exists.

Let $\Q$ be the simplex on $S$ intersected with~\eqref{eq:prefix}, omitting
the moment bound. Let $\Pp$ be the polytope of nonnegative prefix vectors
$(z_1,\ldots,z_a)$ satisfying~\eqref{eq:prefix}, omitting normalization too.
The next section optimizes an arbitrary linear cost on $\Q$.

\section{A normalized prefix oracle}\label{sec:oracle}
\subsection{Backward elimination and canonical schedules}
For any feasible prefix through $i-1$, the next coordinate ranges over
\begin{equation}\label{eq:upper}
 0\le z_i\le U_i:=\frac{ri-\sum_{j<i}(a+j-i)z_j}{a}.
\end{equation}
The numerator is $r>0$ at $i=1$. For $i>1$, the preceding constraint makes
it at least $r+\sum_{j<i}z_j>0$. Setting all future coordinates to zero is
feasible: increasing the horizon decreases the old left-hand side and
increases the right-hand side. Thus eliminating a last coordinate creates
no hidden restriction on an earlier prefix.

For arbitrary costs $c_i$ and a scalar $\alpha$, consider
$\sum_i(c_i-\alpha)z_i$ on $\Pp$. At the last coordinate $i$, the
conditional value is affine in $z_i$ on the entire interval~\eqref{eq:upper}.
If its current coefficient $w_i$ is nonnegative, its minimum is at zero;
if negative, it is at $U_i$. Substituting $U_i$ adds the constant
$w_i ri/a$ and changes each earlier coefficient $w_j$ to
$w_j-w_i(a+j-i)/a$. The remaining value is again affine on every feasible
earlier prefix. Backward induction thus proves equality of the successive
conditional minimum functions, and forward reconstruction attains their
minimum. This proves optimality, as well as feasibility, of the elimination.

Writing $y_i=-\min(w_i,0)/a$ gives the backward recurrence
\begin{equation}\label{eq:y}
 y_i=\max\left\{0,\frac{\alpha-c_i-\sum_{j>i}(a+i-j)y_j}{a}\right\}.
\end{equation}
The accumulated constant is $-r\sum_iiy_i$. Every reduced cost
$c_i-\alpha+\sum_{j\ge i}(a+i-j)y_j$ is nonnegative. A positive $y_i$
selects its tight upper row, and a positive reduced cost selects zero;
therefore the reconstructed prefix attains that dual lower bound.
A zero coefficient permits any point in its interval.

For $B=\{s_1<\cdots<s_k\}\subseteq[a]$, write $v(B)$ for the prefix
schedule supported on $B$ with the prefix constraints tight at those days.
With $s_0=F_0=0$, subtracting consecutive tight constraints gives
\begin{equation}\label{eq:gap}
 v(B)_{s_j}=\frac{s_j-s_{j-1}}a(r+F_{j-1}),\qquad
 F_j=r\left[\prod_{\ell=1}^j
       \left(1+\frac{s_\ell-s_{\ell-1}}a\right)-1\right].
\end{equation}
Here $F_j$ is the accumulated mass. The geometric and gap updates
are those of~\cite[Appendix D.3--D.4]{pmlr-v306-kim26l}.
These schedules are feasible in $\Pp$ but can have mass above one.
Inserting an active day strictly increases the mass: splitting a gap
$u+v$ replaces $1+(u+v)/a$ by
$(1+u/a)(1+v/a)=1+(u+v)/a+uv/a^2$; appending a day adds a factor greater
than one.

\subsection{Nested activation}
Fix arbitrary rational costs $c_t$ on $S$, including signed costs.
Choose a minimizing tail day $t_\star\in S\cap[b,\infty)$, and put
$c_\star=c_{t_\star}$. For $\alpha\le c_\star$,
equation~\eqref{eq:y} gives the dual lower bound
\begin{equation}\label{eq:oracledual}
 \Phi(\alpha)=\alpha-r\sum_{i=1}^a iy_i
\end{equation}
for the normalized optimization on $\Q$. Its slope on an active interval
is $1-m(B)$, where $m(B)=\sum_i v(B)_i$: indeed
$\Phi(\alpha)=\alpha+\min_{z\in\Pp}\sum_i(c_i-\alpha)z_i$.
Increasing $\alpha$ rewards added mass. The optimum occurs when this
slope crosses zero or when the cheapest tail day can fill the missing mass.

For a fixed active set $B$, the active $y_i$ are affine in $\alpha$;
write $\nu_i=dy_i/d\alpha$ there and set $\nu_i=0$ off $B$. Differentiating gives
\begin{equation}\label{eq:speed}
 a\nu_i+\sum_{j>i,\ j\in B}(a+i-j)\nu_j=1\quad(i\in B).
\end{equation}
All active derivatives are positive. More generally, the reduced-slack
speed
\[
 \delta_i=1-\sum_{j>i,\ j\in B}(a+i-j)\nu_j
\]
is positive for every $i$. If there is no later active day, it equals one.
Otherwise let $k$ be the next active day. The equation at $k$ implies
\begin{equation}\label{eq:positive}
 \delta_i=(k-i)\sum_{j\ge k,\ j\in B}\nu_j>0.
\end{equation}
Backward induction proves the assertion, since $\nu_i=\delta_i/a$ at an active
day. Consequently existing positive dual coordinates increase with
$\alpha$, and every inactive reduced slack
\[
 \sigma_i=c_i-\alpha+\sum_{j>i}(a+i-j)y_j
\]
decreases strictly. No day deactivates.

\begin{algorithm}[ht]
\caption{Normalized prefix sweep: costs $c$ on $S$; output $f,\alpha,y$.}
\label{alg:sweep}
\begin{enumerate}
\item Choose a cheapest tail day $t_\star\ge b$. Set $B=\varnothing$ and
$\alpha=\min(c_{t_\star},c_1,\ldots,c_a)$.
\item Add all zero-slack inactive days to $B$; use~\eqref{eq:y} to update
$y$. Let $v_-,v_+$ be the schedules before and after this event, with
masses $m_-,m_+$.
\item If $m_-<1\le m_+$, return the zero-tail law
\begin{equation}\label{eq:massmix}
 f=(1-\eta)v_-+\eta v_+,\qquad \eta=(1-m_-)/(m_+-m_-).
\end{equation}
\item If $\alpha=c_{t_\star}$, fill $1-m_+$ at $t_\star$ and return.
Otherwise advance to the smaller of $c_{t_\star}$ and the next inactive
root $\alpha+\sigma_i/\delta_i$, then repeat from step 2.
\end{enumerate}
\end{algorithm}
At activation, all added dual coordinates are zero, so the old dual
coordinates and current slacks remain unchanged. Activating simultaneous
zero slacks together fixes the convention for ties, including an activation
at the tail cap. The crossing test precedes tail filling.

\begin{lemma}\label{lem:oracle}
The sweep returns a minimizer of $\sum_{t\in S}c_tf_t$ on $\Q$,
attaining~\eqref{eq:oracledual}, with at most one tail day. It takes
$O(b^2+|S|)$ rational operations.
\end{lemma}
\begin{proof}
At a mass crossing both $v_-$ and $v_+$ minimize the unnormalized prefix
problem at the activation price. Every positive old $y_i$ corresponds to
a tight constraint in both schedules, and new $y_i$ equal zero. Their
mixture~\eqref{eq:massmix} is prefix feasible, has mass one, and attains
the normalized dual bound. When the tail cap is reached first, $v(B)$
attains the prefix bound and the remaining mass has cost exactly $\alpha$,
again giving equality. Positivity of the speeds and strict mass increase establish termination.

For fixed $B$, suffix sums of $y_i,iy_i,\nu_i,i\nu_i$ compute all dual
coordinates, slacks and speeds in $O(b)$ operations.
Equation~\eqref{eq:gap} reconstructs a schedule in another $O(b)$ pass.
Scanning all slack roots takes $O(b)$ operations. There are at most $a$
activations and one terminal scan. Selecting the cheapest tail takes
$O(|S|)$ operations. This proves the bound.
\end{proof}

\section{Recovering the exact moment price}\label{sec:price}
Apply Lemma~\ref{lem:oracle} at nonnegative price $\lambda$ to
\begin{equation}\label{eq:q}
 c_t=g(t)+\lambda(t-1),\qquad
 q(\lambda)=\min_{f\in\Q}\{J_p(f)+\lambda\mu(f)\}.
\end{equation}
If an output at zero has moment at most $K$, it solves the robust problem.
Otherwise we maintain oracle outputs $f_-,f_+$ at prices
$\lambda_-<\lambda_+$ with $\mu(f_-)>K\ge\mu(f_+)$.
Each policy defines a line with slope $\mu(f)$ in~\eqref{eq:q}. Their lower
envelope is concave, so its slope decreases as the price increases. We
locate the change across $K$, then mix the two neighboring optimal policies.

Figure~\ref{fig:envelope} illustrates this local transition.

\begin{figure}[ht]
\centering
\begin{tikzpicture}[x=.85cm,y=.65cm,>=Stealth,
 every node/.style={font=\normalsize}]
\draw[->] (0,.1)--(6.6,.1) node[right] {$\lambda$};
\draw[->] (0,.1)--(0,5.3) node[above] {$q(\lambda)$};
\draw[densely dashed,gray] (0,.7)--(6,4.9);
\draw[densely dashed,gray] (0,1.9)--(6,3.7);
\draw[very thick,blue] (0,.7)--(3,2.8)--(6,3.7);
\draw[dashed] (2.1,.1)--(2.1,2.17);
\draw[dashed] (3.9,.1)--(3.9,3.07);
\draw[dotted] (3,.1)--(3,2.8);
\fill[blue] (3,2.8) circle[radius=2pt];
\node[below] at (2.1,.1) {$\lambda_-$};
\node[below] at (3,.1) {$\kappa$};
\node[below] at (3.9,.1) {$\lambda_+$};
\node[anchor=west,align=left] at (7,4.1)
 {$\ell_-(\lambda)=J_p(f_-)+\lambda\mu(f_-)$\\$\mu(f_-)>K$};
\node[anchor=west,align=left] at (7,2.2)
 {$\ell_+(\lambda)=J_p(f_+)+\lambda\mu(f_+)$\\$\mu(f_+)\le K$};
\end{tikzpicture}
\caption{Local schematic around the isolated price $\kappa$. The blue lower
envelope changes from the higher-moment policy to the lower-moment policy;
both are optimal at the intersection. The dashed verticals show the bracket.}
\Description{Two affine policy lines intersect at kappa. Their concave lower
envelope has a larger slope on the left than on the right. Bracket endpoints
lambda minus and lambda plus lie on the two sides of the intersection.}
\label{fig:envelope}
\end{figure}

\begin{lemma}\label{lem:isolation}
With $Z$ from~\eqref{eq:bound}, price $Z+1$ supplies a right endpoint.
After $O(\log Z)$ bisection steps the bracket contains a common optimal
price for both stored policies, recoverable exactly from their costs and
moments.
\end{lemma}
\begin{proof}
A vertex of $\Q$ has at most $b$ positive coordinates. Its coordinates
are obtained from normalization and independent tight prefix rows, with
integer coefficients of magnitude at most $b$. Their determinant has
absolute value at most $D=b!b^b$. Cramer's rule gives a common coordinate
denominator $q_R|\det|$. Its objective and moment therefore each have
denominator at most $H$, and lie in $[0,L]$.

The intersection price of two vertex lines $J+\lambda\mu$ with distinct
moments has numerator and denominator at most $Z=2LH^4$ in reduced form:
each difference has numerator at most $2LH^2$ and denominator at most
$H^2$. Hence all positive intersections are at most $Z$, and distinct
nonnegative intersections, including zero, are separated by at least
$1/Z^2$. Beyond $Z$, every minimizing vertex has minimum possible moment;
so does every convex combination of them. Robust feasibility gives that
this minimum moment is at most $K$. Thus $\lambda_+=Z+1$ is valid.

At each bisection midpoint, replace the left endpoint if the returned
moment exceeds $K$, and otherwise replace the right endpoint. Adding the
two optimality inequalities at distinct prices proves that oracle moments
are nonincreasing in price, without assuming vertex outputs or unique
optima. After at most
$\lceil\log_2((Z+1)Z^2)\rceil+1$ iterations the bracket width is strictly
less than $1/Z^2$.

The function $q$ is the lower envelope of finitely many vertex lines.
The strict moment bracket forces a breakpoint in the closed interval:
without one, all minimizing lines have the same moment throughout it.
There is at most one distinct intersection price in that interval.
On each side of this price all minimizing vertex lines have the same
objective and moment. By continuity, both stored endpoint policies are
optimal at the breakpoint $\kappa$. If an endpoint itself is the breakpoint,
its possibly nonvertex output is already optimal there. 
\end{proof}

The two stored moments differ, so their line intersection and a feasible
mixture are
\begin{equation}\label{eq:recover}
 \kappa=\frac{J_p(f_+)-J_p(f_-)}{\mu(f_-)-\mu(f_+)},\qquad
 \theta=\frac{K-\mu(f_+)}{\mu(f_-)-\mu(f_+)},\qquad
 f=\theta f_-+(1-\theta)f_+.
\end{equation}
Here $0\le\theta<1$, $\mu(f)=K$, and both endpoints minimize
$q(\kappa)$, so their mixture does too. Recompute the prefix oracle dual
at $\kappa$ to obtain $\alpha,y$.

\begin{algorithm}[ht]
\caption{Exact moment-price recovery; $Z$ is defined in~\eqref{eq:bound}.}
\label{alg:price}
\begin{enumerate}
\item Run Algorithm~\ref{alg:sweep} at $\lambda=0$. If $\mu\le K$, return
this policy and its dual. Otherwise store it as $f_-$ at $\lambda_-=0$.
\item Store the output at $\lambda_+=Z+1$ as $f_+$.
While $\lambda_+-\lambda_-\ge1/Z^2$, query the midpoint price.
If its moment exceeds $K$, replace the left endpoint; otherwise replace
the right endpoint.
\item Recover $\kappa,\theta,f$ by~\eqref{eq:recover}. Run the sweep at
$\kappa$ to obtain its dual and return $f,\kappa,\alpha,y$.
\end{enumerate}
\end{algorithm}

\subsection{Global certificates and complexity}
The oracle dual at $\kappa$ gives~\eqref{eq:certificate} on $S$.
Equation~\eqref{eq:dominance} extends it to every omitted $t\ge b$;
all $t<b$ already belong to $S$. For any robust law, taking expectations
of~\eqref{eq:certificate} gives
\[
 J_p(f)\ge\alpha-r\sum_{x=1}^{b-1}xy_x-\lambda K.
\]
The expectation is well defined because $g$ is bounded, the prefix sum
is finite, and robustness gives finite moment.
The zero-price feasible output or the mixture~\eqref{eq:recover} attains
equality. Each endpoint oracle uses at most one tail day. Their mixture
therefore uses at most two, proving the certificate and support assertions
of Theorem~\ref{thm:main}.

Let $\ell$ be the binary encoding length of the input, with a sparse list
of forecast atoms. Then $\log H=O(\ell+b\log b)$ and
$\log Z=O(\ell+b\log b)$. Sorting forecast atoms and using prefix
probability and first-moment sums computes $g$ on $S$ in
$O(|S|+|\supp p|\log|\supp p|)$ rational operations.
The simple reference program directly sums the costs and remains
polynomial. Excluding preprocessing, the total number of rational operations is
\[
 O\bigl((b^2+|S|)(\ell+b\log b)\bigr).
\]

To bound intermediate bits, fix an active set of size $k\le a$, and let
$Q_c$ be a common denominator of the original costs in this oracle call,
$M=\max(1,\max_t|c_t|)$. The triangular solve gives
$y_i=u_i\alpha+v_i$, with denominators dividing $a^k$ and $Q_ca^k$,
respectively. Each substitution divides once by $a$, and each normalized
off-diagonal coefficient has magnitude at most one. Induction on the
remaining sum of absolute coefficients gives $|u_i|\le2^k$ and
$|v_i|\le2^kM$. Every inactive slack $U\alpha+V$ therefore has the
same common denominators and bounds
\[
 |U|\le1+ak2^k,\qquad |V|\le M(1+ak2^k).
\]
Its root has $O(\log Q_c+\log M+b\log b)$ bits. At each event this
system is expressed in the original costs, so successive roots do not
recursively enlarge their denominators. Evaluation and suffix sums add
only polynomially many bits. Canonical schedules have common denominator
dividing $q_Ra^a$ by~\eqref{eq:gap}, and numerators of
$O(\operatorname{len}(R)+b\log b)$ bits; normalization uses one ratio
of mass differences.

After $k$ bisections, the price denominator is at most $2^k$ and the
numerator at most $(Z+1)2^k$, with $k=O(\log Z)$. Adjusted costs have
common denominator dividing $q_p\operatorname{den}(\lambda)$ and
polynomial bit length in $b,\ell$. The preceding bounds apply to every
oracle call. Finally~\eqref{eq:recover} and the dual reconstruction use
polynomial-length rationals. Standard exact rational arithmetic therefore
proves polynomial total bit time, completing Theorem~\ref{thm:main}.

\section{The uniform-forecast frontier}\label{sec:uniform}
For a uniform forecast on $\{1,\ldots,b\}$, every day after the forecast
support has the same cost. The moment cap determines how much buying must
be moved into the forecast interval; the prefix constraints limit that
mass. Filling successive geometric prefix coordinates gives a complete
frontier, and strict dual prices determine the policy uniquely.

Put $m=(b+1)/2$. Define $z_i,d_i$ for $1\le i\le b$,
$A_j,C_j,r_j$ for $0\le j\le b$, and $\lambda_j,s_j$ for
$1\le j\le b$ by
\begin{align}
 z_i&=\frac{b^{i-1}}{a^i},&
 A_j&=(b/a)^j-1,&
 C_j&=(2b-j)(b/a)^j-2b,&
 r_j&=\frac b{b+C_j},\label{eq:uniformdata}\\
 d_i&=\frac{(b+1-i)(b+i-2)}{2b},&
 \lambda_j&=\frac{b+j-2}{2b},&
 s_j&=\sum_{i<j}d_i z_i-\lambda_j(b+C_{j-1}).\label{eq:uniformslope}
\end{align}
In particular, $A_0=C_0=0$.

\begin{theorem}\label{thm:uniform}
For the uniform forecast, $1=r_0>\cdots>r_b=\rho_b-1$.
If $r_j\le r\le r_{j-1}$, the unique optimal law is
\begin{equation}\label{eq:uniformpolicy}
 f_i=rz_i\ (i<j),\qquad
 f_j=\frac{b-r(b+C_{j-1})}{b+1-j},\qquad
 f_{b+1}=1-\sum_{i\le j}f_i,
\end{equation}
with all other masses zero. Adjacent formulas agree at shared endpoints.
Its value on segment $j$ is $m+\lambda_j b+s_jr$.
For $R\ge2$, day $b+1$ alone is optimal with value $m$.
The value has exactly $b$ strictly decreasing affine pieces followed by
this plateau. The shared frontier representation, a single policy/value
query, and a dual query each take $O(b)$ rational operations; every rational
uses $O(b\log b+\operatorname{len}(R))$ bits.
For $r_b<r<r_{b-1}$, all $b+1$ masses on $\{1,\ldots,b+1\}$ are positive.
Thus both support bounds of Theorem~\ref{thm:main} are simultaneously sharp.
\end{theorem}

The capped-geometric law in~\eqref{eq:uniformpolicy} agrees with the known
one-hot-$b$ optimizer in the range described in Section~\ref{sec:intro}.
The plateau's unconstrained forecast behavior also appears in the discrete
Bayesian model of Kang, Park and Fan~\cite{kang2025learningaugmentedskirentaldiscrete}.
We prove the frontier and uniqueness using its forecast-specific dual.

\begin{proof}
Direct summation gives $g(i)=m+d_i$ for $i\le b$, and $g(t)=m$ for
$t\ge b+1$. Geometric summation, or subtraction of consecutive expressions,
gives
\[
 A_j=\sum_{i\le j}z_i,\quad
 C_j=\sum_{i\le j}(b+1-i)z_i,\quad
 M_j:=\sum_{i\le j}(i-1)z_i=b-(b-j)(b/a)^j.
\]
Thus $C_j$ increases strictly, $r_b=a^b/(b^b-a^b)$, and $M_j\le b$,
strictly for $j<b$. At $r=r_j$, the policy has $f_i=r_jz_i$ for
$i\le j$; at $r=r_{j-1}$, $f_j=0$. These endpoint identities prove
agreement and, by affinity, $0\le f_j\le rz_j$ throughout the interval.

The full geometric prefix makes its rows tight, since
$af_i=r+\sum_{k<i}f_k$. The displayed policy has the first $j-1$ rows
tight and the $j$th row feasible when $j<b$. With no further early mass,
subsequent left sides decrease while their right sides increase, so all
prefix rows hold. At endpoint $r_j$ the residual tail is $1-r_jA_j\ge0$:
using $C_j=bA_j-M_j$, this inequality is equivalent to $M_j\le b$.
The other endpoint is the preceding policy, so affinity proves tail
nonnegativity on the whole segment. Finally,
\begin{equation}\label{eq:uniformmoment}
 \mu=b-\sum_{i\le j}(b+1-i)f_i=rb.
\end{equation}
This proves feasibility, including the classical endpoint where the tail
vanishes. For $r\ge1$, day $b+1$ has moment $b\le rb$ and cost $m$.

For segment $j$, choose $\lambda=\lambda_j>0$,
$\alpha=m+\lambda b$, and $y_i=0$ for $i\ge j$. Choose $y_i$ backward
for $i<j$ to make~\eqref{eq:certificate} an equality at those days.
It is already an equality at $j$, since $d_j=\lambda(b+1-j)$.
For $j\ge2$, the last early price satisfies
$ay_{j-1}=(b-j+2)/(2b)>0$. Subtracting consecutive equality equations gives
\begin{equation}\label{eq:uniformdual}
 ay_i-\sum_{k=i+1}^{j-1}y_k=\frac{b+j-2i}{2b}>0\qquad(1\le i<j-1).
\end{equation}
Hence all $y_i$ for $i<j$ are positive. For $j<t\le b$, the reduced
cost is $(b+1-t)(t-j)/(2b)>0$; at $b+1$ it is zero, and for $t>b+1$
it is $\lambda(t-b-1)>0$. This certifies the inequality for every buy-day.
The displayed policy uses only equality days, tightens every positively
priced row, and satisfies~\eqref{eq:uniformmoment}, proving optimality.

Conversely, equality in the dual bound forces every optimum to have support
in $\{1,\ldots,j,b+1\}$, its first $j-1$ prefix rows tight, and moment
$rb$. These triangular equations and normalization uniquely give
\eqref{eq:uniformpolicy}, including shared endpoints and $r=1$.
Substitution into $J_p=m+\sum_i d_if_i$ gives the value claimed. Its slopes
satisfy
\[
 s_1=-(b-1)/2<0,\qquad s_{j+1}-s_j=-\frac{b+C_j}{2b}<0.
\]
As $r$ increases the segments are traversed in reverse $j$ order, so all
$b$ distinct negative slopes occur, followed by zero. In the lowest open
segment ($j=b$), the first $b$ masses are positive. The tail is affine,
zero at $r_b$ and positive at $r_{b-1}$, so it too is positive inside.
Uniqueness excludes every smaller-support optimum, even on later days.

Compute successive $z_i$ and accumulate $A,C,\sum d_iz_i$ to store one
slope, intercept and breakpoint per $j$. This takes $O(b)$ operations
and shares the prefix data across segments. An interval scan, policy
construction and backward dual solve each take $O(b)$ operations.
Powers $b^j,a^j$ have $O(b\log b)$ bits. The sums share denominator
$a^b$, the $d_i$ have denominator $2b$, and the dual recurrence has depth
at most $b$ with diagonal $a$. All resulting rationals have
$O(b\log b+\operatorname{len}(R))$ bits. This bound concerns the compact
representation and one query, rather than printing $b$ dense policy vectors.
\end{proof}

\begin{figure}[ht]
\centering
\begin{tikzpicture}[x=10cm,y=2.5cm,>=Stealth,every node/.style={font=\normalsize}]
\draw[->] (1.38,1.85)--(2.13,1.85) node[right] {$R$};
\draw[->] (1.38,1.85)--(1.38,3.02) node[above] {$J^*(R)$};
\draw[thick,blue] (1.421052632,2.842105263)--(1.5,2.625)--(1.666666667,2.333333333)--(2,2)--(2.1,2);
\foreach \x/\y in {1.421052632/2.842105263,1.5/2.625,1.666666667/2.333333333,2/2}
  \fill[blue] (\x,\y) circle[radius=1.7pt];
\node[below left] at (1.421052632,1.85) {$27/19$};
\node[below right] at (1.5,1.85) {$3/2$};
\node[below] at (1.666666667,1.85) {$5/3$};
\node[below] at (2,1.85) {$2$};
\node[above right] at (1.421052632,2.842105263) {$54/19$};
\node[above right] at (1.5,2.625) {$21/8$};
\node[above right] at (1.666666667,2.333333333) {$7/3$};
\node[above] at (2.045,2) {$2$};
\end{tikzpicture}
\caption{The exact uniform frontier at $b=3$: three affine pieces and the
plateau. On $27/19<R<3/2$, the unique optimum needs all four buy-days.}
\Description{For purchase cost three and a uniform forecast, the value decreases through the four exact endpoint values shown, then stays constant at two.}
\label{fig:frontier}
\end{figure}

\paragraph{Further questions.}
Further directions include worst-case bounds on frontier size, direct
structural formulas for arbitrary forecasts, and improved dependence on
numerical $b$. The general-forecast structural frontier is left for future work.

\paragraph{Computational and AI assistance.}
Codex assisted with proof development, exact checks, source comparison and
drafting. A separate model-context pass checked the proofs and certificates.

\bibliographystyle{alpha}
\bibliography{references}
\clearpage\onehalfspacing
\end{document}
